have the Chow ring @@M0@@ of cycle classes modulo rational equivalence [4]. As variants of this ring, which is moreover defined under substantially more general considerations, there are the rings @@M1@@ (the associated graded of the @@M2@@-filtration of @@M3@@) and @@M4@@ (the associated graded of the filtration of @@M5@@ by the codimension of supports). There are theories of Chern classes with values in each of these three rings (cf. [9] for @@M6@@, and this seminar for @@M7@@, which maps to @@M8@@). There is also a cohomological theory (@@M9@@-adic, say, with @@M10@@ prime to the characteristic of @@M11@@) of Chern classes (SGA 5 VII):
@@M12@@
(where @@M13@@ is a vector bundle on @@M14@@, or indeed any element of @@M15@@). These cohomological theories are related by the canonical homomorphisms (compatible with the multiplicative structures)
\sgaStableTarget{sga6.e14.eq.4-1}@@M16@@
where we have set
@@M17@@
The homomorphism @@M18@@ comes from the fact that the @@M19@@-filtration of @@M20@@ is finer than the topological filtration by the codimension of the support (Exp. X), and it is an isomorphism modulo torsion (Exp. VII 4.11). The homomorphism @@M21@@ is defined, for example, in [RRR], and associates to the class of the cycle of codimension @@M22@@ associated with a closed subscheme @@M23@@ of @@M24@@ of codimension @@M25@@ the element @@M26@@ of @@M27@@ mod @@M28@@. It is a surjective homomorphism, which we shall also see (4.2) to be an isomorphism modulo torsion. Thus, \uline{modulo torsion the three rings} @@M29@@, @@M30@@ \uline{and} @@M31@@ \uline{are canonically isomorphic}, so \uline{a posteriori} the Riemann--Roch formulas expressed using one or another of these rings (tensored with @@M32@@) are equivalent. The graded ring @@M33@@, used particularly in this seminar as a substitute for the Chow ring, has the appreciable advantage of being defined for every locally ringed topos and, in particular, for every scheme
